\begin{center}
		
	Министерство науки и высшего образования Российской Федерации \\
	ФЕДЕРАЛЬНОЕ ГОСУДАРСТВЕННОЕ АВТОНОМНОЕ \\
	ОБРАЗОВАТЕЛЬНОЕ УЧРЕЖДЕНИЕ ВЫСШЕГО ОБРАЗОВАНИЯ \\
	«НАЦИОНАЛЬНЫЙ ИССЛЕДОВАТЕЛЬСКИЙ ЯДЕРНЫЙ УНИВЕРСИТЕТ	«МИФИ» \\
	
	\vspace{0.5cm}
	
	Институт физико-технических интеллектуальных систем \\
	
	\vspace{0.5cm}
	
	Кафедра №18 «Конструирование приборов и установок» \\
	
	\vspace{2cm}
	
	\normalsize
	ОТЧЕТ \\
	\vspace{0.1em}
	ПО ПРОИЗВОДСТВЕННОЙ ПРАКТИКЕ
	
	\vspace{1cm}
	
	по теме: \\
	\vspace{0.1em}
	РАЗРАБОТКА СИСТЕМЫ КОНТРОЛЯ \\ ГЕОМЕТРИЧЕСКИХ ПАРАМЕТРОВ ЖЕЛЕЗНОДОРОЖНОГО ПОЛОТНА \\ НА ОСНОВЕ МИКРОМЕХАНИЧЕСКИХ ДАТЧИКОВ \\
	
\end{center}

\vspace{3cm}

% Блок с подписями	
\begin{table}[h!]
	\centering
	\begin{tabular}{l @{\hspace{15pt}} l @{\hspace{15pt}} r}
		Исполнитель: 				& 					& \\
		Студент группы Б22-604 		& \line(1,0){90} 	& Р.И. Романовский \\
		& & \\
		Руководитель практической 	& & \\
		подготовки от НИЯУ МИФИ, 	& 					& \\
		старший преподаватель каф. №18	& \line(1,0){90} 	& А.И. Максимкин \\
		& & \\
		Руководитель практической 	& & \\
		подготовки от ЗАО \enquote{ПИК Прогресс}, & & \\
		начальник отдела			& \line(1,0){90} 	& М.А. Глазков \\
		& & \\
		Зав. каф. №18, к.т.н. 	& \line(1,0){90} 	& Р.Е. Невский \\
	\end{tabular}
\end{table}

\vspace{1cm}

\begin{center}
	Москва 2026
\end{center}

\newpage